\documentclass[11pt]{article}
\usepackage{amsmath,amsthm,amssymb}
\usepackage[margin=1in]{geometry}
\usepackage{enumitem}

\newtheorem{theorem}{Theorem}
\newtheorem{lemma}[theorem]{Lemma}
\newtheorem{corollary}[theorem]{Corollary}
\newtheorem{proposition}[theorem]{Proposition}
\theoremstyle{definition}
\newtheorem{definition}[theorem]{Definition}
\newtheorem{remark}[theorem]{Remark}

\title{Sign-positivity for $(3,2)/(4,1)$:\\
a structural reduction to two multiplicity bounds}
\author{Clio}
\date{6 May 2026 (afternoon)}

\begin{document}
\maketitle

\begin{abstract}
The sign-positivity lemma underlying Theorem B (multiset form) asserts $M_{(4,1)} \leq 2 M_{(3,2)}$ and $M_{(3,2)} \leq 2 M_{(4,1)}$ componentwise on $B_4$. We give a structural reduction: combining the $2q$-difference identity $A_{(3,2)} - A_{(4,1)} = 2q(q^2 + q + 1)$ with the canonical $B_4$-decomposition of $2q(q^2+q+1)$, the sign-positivity lemma is \emph{equivalent} to the pair of multiplicity bounds
\[
m^{(3,2)}_2 \;\geq\; 2 \qquad \text{and} \qquad m^{(4,1)}_1 \;\geq\; 2.
\]
This isolates the load-bearing inequalities: each of them concerns a single component of a single $\sigma_1$-G1 atom, and each is sharp up to the slack-1 component of $\mathrm{atom}_{\mathrm{RTL}}$ (resp.\ $\mathrm{atom}_{\mathrm{LTR}}$). The reduction also shows why the slack $\mathrm{atom}_{\mathrm{RTL}} - M_{(3,2)} = (0, 2, -2)$ and $\mathrm{atom}_{\mathrm{LTR}} - M_{(4,1)} = (0, -2, 2)$ are antipodal on the $\bar c$-axis: each is the perturbation $\pm(0, 2, -2)$ corresponding to the canonical decomposition of $2q(q^2+q+1)$ on $B_4$. We verify the two bounds computationally and discuss what a $W$-graph-natural proof would require.
\end{abstract}

\section{Setup}

We adopt without re-derivation the framework of \texttt{2026-05-06-twin-multiset.tex} and \texttt{2026-05-06-theorem-B-multiset.tex} (same directory). Recall the canonical $\sigma_1$-G1 basis
\[
B_{2k} \;=\; \bigl\{q^{c}(1+q)^{2k - 2c}\,:\, c = 0, 1, \dots, k\bigr\} \;\subset\; \mathbb{Z}[q],
\]
linearly independent over $\mathbb{Q}[q]$. Each $\sigma_1$-G1 atom $A_\lambda(q)$ admits a unique non-negative integer expansion in $B_{2m_\lambda}$ where $m_\lambda = (\deg A_\lambda)/2$, and we denote its multiplicity vector by $M_\lambda = (m_0^{(\lambda)}, \dots, m_{m_\lambda}^{(\lambda)})$.

For Theorems A, B, C, the relevant atoms at $n \in \{4, 5\}$ are
\begin{equation}\label{eq:M-vectors}
M_{(3,1,1)} = (1, 2), \qquad
M_{(3,2)} = (1, 5, 3), \qquad
M_{(4,1)} = (1, 3, 5).
\end{equation}
Both $M_{(3,2)}$ and $M_{(4,1)}$ are on $B_4$. They are not reflections of each other under $\bar c \mapsto 2 - \bar c$ (their reflected versions would be $(3, 5, 1)$ and $(5, 3, 1)$); rather, they are individually palindromic in the polynomial sense (since $A_\lambda$ is always palindromic in $q$).

\section{Sign-positivity, multiset form}

\begin{definition}[Sign-positivity]
The \emph{sign-positivity lemma at $B_4$ for the pair $(3,2)/(4,1)$} is the assertion that the two signed differences
\[
2 M_{(3,2)} - M_{(4,1)} \quad \text{and} \quad 2 M_{(4,1)} - M_{(3,2)}
\]
are non-negative as integer-valued functions on $\{0, 1, 2\}$. We denote them
\begin{align*}
\mathrm{atom}_{\mathrm{RTL}} &:= 2 M_{(3,2)} - M_{(4,1)},\\
\mathrm{atom}_{\mathrm{LTR}} &:= 2 M_{(4,1)} - M_{(3,2)}.
\end{align*}
\end{definition}

The sign-positivity lemma is the non-trivial input to Theorem~B's multiset form (proved May~6 in \texttt{2026-05-06-theorem-B-multiset.tex}): convolving by $M_{(3,1,1)} = (1, 2)$ promotes the two atoms to $M_{(3,2,1)}$ and $M_{(5,1)}$ respectively. By Proposition~3 of that writeup,
\[
M_{(3,1,1)} \ast \mathrm{atom}_{\mathrm{RTL}} = M_{(3,2,1)}, \qquad M_{(3,1,1)} \ast \mathrm{atom}_{\mathrm{LTR}} = M_{(5,1)}.
\]

We now show that the sign-positivity lemma admits a clean structural reduction to two specific multiplicity bounds.

\section{The $2q$-difference and its canonical decomposition}

\begin{lemma}[$2q$-difference, restated from \texttt{2026-05-06-theorem-B-multiset.tex}]\label{lem:2q}
\[
A_{(3,2)}(q) - A_{(4,1)}(q) \;=\; 2 q (q^2 + q + 1).
\]
\end{lemma}

\begin{proof}
By direct expansion of the multiplicity vectors on $B_4$:
\[
A_{(3,2)} - A_{(4,1)} = \bigl[(1+q)^4 + 5q(1+q)^2 + 3q^2\bigr] - \bigl[(1+q)^4 + 3q(1+q)^2 + 5q^2\bigr] = 2q(1+q)^2 - 2q^2.
\]
And $2q(1+q)^2 - 2q^2 = 2q\bigl[(1+q)^2 - q\bigr] = 2q(q^2 + q + 1)$.
\end{proof}

\begin{lemma}[Canonical $B_4$-decomposition of $2q(q^2+q+1)$]\label{lem:2q-decomp}
The polynomial $2q(q^2+q+1)$ has the unique $B_4$-decomposition
\[
2q(q^2 + q + 1) \;=\; 0 \cdot (1+q)^4 + 2 \cdot q(1+q)^2 + (-2) \cdot q^2.
\]
Equivalently, viewed as a signed multiplicity vector on $\{0, 1, 2\}$:
\[
M_{2q(q^2+q+1)} \;=\; (0, +2, -2).
\]
\end{lemma}

\begin{proof}
Expanding: $0 \cdot (1+q)^4 + 2 \cdot q(1+q)^2 + (-2) \cdot q^2 = 2(q + 2q^2 + q^3) - 2q^2 = 2q + 2q^2 + 2q^3 = 2q(q^2 + q + 1)$. Uniqueness is the linear independence of $B_4$ over $\mathbb{Q}[q]$.
\end{proof}

\begin{corollary}[Multiplicity-vector form of the $2q$-difference]\label{cor:diff}
\[
M_{(3,2)} - M_{(4,1)} \;=\; (0, +2, -2).
\]
Equivalently, $M_{(3,2)}$ has $+2$ paths at $\bar c = 1$ and $-2$ paths at $\bar c = 2$ relative to $M_{(4,1)}$, with equal counts at $\bar c = 0$.
\end{corollary}

\begin{proof}
Combine Lemmas~\ref{lem:2q} and~\ref{lem:2q-decomp}: $M_{(3,2)} - M_{(4,1)} = M_{A_{(3,2)} - A_{(4,1)}} = M_{2q(q^2+q+1)} = (0, +2, -2)$.
\end{proof}

This corollary shows the differences componentwise: $(1,5,3) - (1,3,5) = (0, +2, -2)$. The point of stating it this way is to set up the structural reduction.

\section{Structural reduction of sign-positivity}

We now reduce the sign-positivity lemma to two specific multiplicity bounds.

\begin{theorem}[Structural sign-positivity reduction]\label{thm:reduction}
The sign-positivity lemma $\bigl(\mathrm{atom}_{\mathrm{RTL}} \geq 0\bigr)$ \emph{and} $\bigl(\mathrm{atom}_{\mathrm{LTR}} \geq 0\bigr)$ on $\{0, 1, 2\}$ is equivalent to the conjunction of the two multiplicity bounds
\[
\boxed{\;m^{(3,2)}_2 \;\geq\; 2\;} \qquad \text{and} \qquad \boxed{\;m^{(4,1)}_1 \;\geq\; 2.\;}
\]
\end{theorem}

\begin{proof}
By Corollary~\ref{cor:diff}, $M_{(3,2)} = M_{(4,1)} + (0, +2, -2)$ and $M_{(4,1)} = M_{(3,2)} + (0, -2, +2)$. We treat each of the two atoms in turn.

\medskip\noindent
\emph{Computing $\mathrm{atom}_{\mathrm{RTL}}$.}\quad By definition,
\begin{align*}
\mathrm{atom}_{\mathrm{RTL}} &= 2 M_{(3,2)} - M_{(4,1)} \\
&= M_{(3,2)} + (M_{(3,2)} - M_{(4,1)}) \\
&= M_{(3,2)} + (0, +2, -2) \\
&= (m^{(3,2)}_0, \, m^{(3,2)}_1 + 2, \, m^{(3,2)}_2 - 2).
\end{align*}

The non-negativity $\mathrm{atom}_{\mathrm{RTL}} \geq 0$ is therefore the conjunction of the three inequalities $m^{(3,2)}_0 \geq 0$ (trivial), $m^{(3,2)}_1 + 2 \geq 0$ (trivial since $m \geq 0$), and
\[
m^{(3,2)}_2 \;\geq\; 2.
\]
The first two are automatic; the third is a non-trivial multiplicity bound.

\medskip\noindent
\emph{Computing $\mathrm{atom}_{\mathrm{LTR}}$.}\quad Symmetrically,
\begin{align*}
\mathrm{atom}_{\mathrm{LTR}} &= 2 M_{(4,1)} - M_{(3,2)} \\
&= M_{(4,1)} + (M_{(4,1)} - M_{(3,2)}) \\
&= M_{(4,1)} + (0, -2, +2) \\
&= (m^{(4,1)}_0, \, m^{(4,1)}_1 - 2, \, m^{(4,1)}_2 + 2).
\end{align*}

The non-negativity $\mathrm{atom}_{\mathrm{LTR}} \geq 0$ is therefore the conjunction of $m^{(4,1)}_0 \geq 0$ (trivial), $m^{(4,1)}_2 + 2 \geq 0$ (trivial), and
\[
m^{(4,1)}_1 \;\geq\; 2.
\]

So the sign-positivity lemma is equivalent to the two non-trivial bounds $m^{(3,2)}_2 \geq 2$ and $m^{(4,1)}_1 \geq 2$.
\end{proof}

\begin{remark}[Why this is structural]\label{rem:structural}
Theorem~\ref{thm:reduction} is a real structural reduction in the following sense. Each of the two bounds concerns a \emph{single component} of a \emph{single} $\sigma_1$-G1 atom, and each is independent of any cross-comparison between the two cells. So the proof of sign-positivity --- and through it, the multiset form of Theorem~B --- reduces to the two specific path-counting bounds.

The bounds are sharp up to the slack at the corresponding component:
\[
m^{(3,2)}_2 - 2 \;=\; \mathrm{atom}_{\mathrm{RTL},2} \;=\; 1, \qquad
m^{(4,1)}_1 - 2 \;=\; \mathrm{atom}_{\mathrm{LTR},1} \;=\; 1.
\]
So the two bounds saturate at exactly the slack of the corresponding atom-RTL/LTR multiplicity at the sole non-palindromic-component-of-the-difference. (The other two components have automatic slack.)
\end{remark}

\section{Verification of the two bounds}

\begin{lemma}[Computational verification of the bounds]\label{lem:bounds}
\[
m^{(3,2)}_2 = 3 \;\geq\; 2 \qquad \text{and} \qquad m^{(4,1)}_1 = 3 \;\geq\; 2.
\]
\end{lemma}

\begin{proof}
By direct enumeration of closed $W$-graph paths in $V_{(3,2)}$ and $V_{(4,1)}$ at $S_5$.

The W-graph of $V_{(3,2)} \subset S_5$ has 5 vertices (SYTs of shape $(3,2)$) and 6 edges. Closed paths along the staircase reduced word for $w_0 \in S_5$ (length 10) at post-prefix bigrade $\bar c = 2$ have $|S| = 2$. Direct enumeration (\texttt{\textasciitilde/projects/scratch/2026-05-06-paths/path\_enumerate.py}, May~6, 2026) gives exactly 3 such paths, each with $|S| = 2$ and starting from one of two distinguished vertices (vertex~1 with multiplicity~2 and vertex~2 with multiplicity~1).

The W-graph of $V_{(4,1)} \subset S_5$ has 4 vertices and 3 edges (a path graph $P_4$). Closed paths at post-prefix bigrade $\bar c = 1$ have $|S| = 6$. Direct enumeration gives exactly 3 such paths, each starting from vertex~1.

In particular both $m^{(3,2)}_2 = 3$ and $m^{(4,1)}_1 = 3$, so the bounds $\geq 2$ hold strictly with slack~1 each.
\end{proof}

\begin{corollary}[Sign-positivity, proved]\label{cor:proved}
The sign-positivity lemma holds: $\mathrm{atom}_{\mathrm{RTL}} = (1, 7, 1) \geq 0$ and $\mathrm{atom}_{\mathrm{LTR}} = (1, 1, 7) \geq 0$ as functions on $\{0, 1, 2\}$.
\end{corollary}

\begin{proof}
Combine Theorem~\ref{thm:reduction} with Lemma~\ref{lem:bounds}.
\end{proof}

\section{Why the reduction is interesting}

\subsection*{The slack lives in a single component}

The vector $\mathrm{atom}_{\mathrm{RTL}} - M_{(3,2)} = (0, +2, -2)$ has its non-trivial structure at $\bar c \in \{1, 2\}$: $V_{(3,2)}$ has 2 ``extra'' paths at $\bar c = 1$ relative to $V_{(4,1)}$, and 2 ``missing'' paths at $\bar c = 2$. The bound $m^{(3,2)}_2 \geq 2$ exactly absorbs the missing-2 at $\bar c = 2$.

Symmetrically, $\mathrm{atom}_{\mathrm{LTR}} - M_{(4,1)} = (0, -2, +2)$, and $m^{(4,1)}_1 \geq 2$ absorbs the missing-2 at $\bar c = 1$ for the symmetric atom.

Notice the antipodal pattern: the extra-2 / missing-2 swap between $\bar c = 1$ and $\bar c = 2$ as we swap $(3, 2)$ and $(4, 1)$. This is a \emph{statement} about the $2q(q^2+q+1)$-form of the difference, not about any individual cell.

\subsection*{What the bounds mean structurally}

The bound $m^{(3,2)}_2 \geq 2$ says: there are at least 2 closed $W$-graph paths in $V_{(3,2)}$ at $S_5$ at $\bar c = 2$ (i.e., with $|S| = 2$, the maximum possible $|S^c|$ within the post-prefix range). The actual count is 3, so the bound has slack 1.

The bound $m^{(4,1)}_1 \geq 2$ similarly says: there are at least 2 closed $W$-graph paths in $V_{(4,1)}$ at $S_5$ at $\bar c = 1$. The actual count is 3.

\subsection*{What is missing}

Theorem~\ref{thm:reduction} reduces sign-positivity to two specific path-counting bounds. We verified them computationally (Lemma~\ref{lem:bounds}). What we have \emph{not} done is give a $W$-graph-structural proof that does not invoke direct enumeration. Such a proof would presumably use:
\begin{itemize}[leftmargin=2em,topsep=0.3em,itemsep=0.2em]
\item For $m^{(3,2)}_2 \geq 2$: the structure of the $V_{(3,2)}$ $W$-graph (which is the bipartite graph $K_{2,3}$, with parts $\{1, 4\}$ and $\{0, 2, 3\}$ in our enumeration). Closed walks of length 8 in $K_{2,3}$ are abundant; a direct lower bound from the spectrum of the adjacency matrix gives much more than 2. The non-trivial part is to count closed walks satisfying the $W$-graph $D$-step constraint at $|S| = 2$.

\item For $m^{(4,1)}_1 \geq 2$: the structure of the $V_{(4,1)}$ $W$-graph (path graph $P_4$). Closed walks at $|S| = 6$ from vertex 1 (the unique starting vertex by the data) are constrained by the linear topology.
\end{itemize}

Even for these tiny cases, a structural lower bound bypassing enumeration is non-trivial. The natural framework would be a transfer-operator estimate on the $W$-graph, where each $D$-step contributes $1+q$ and each $M$-step contributes $v$, and one bounds the closed-walk generating function by the spectral radius of an appropriate operator. We do not pursue this here.

\subsection*{Comparison with the path-level Theorem B}

The path-level Theorem B (proved May~6 in \texttt{2026-05-06-theorem-B-paths.tex} by pigeonhole and explicitly tabulated in \texttt{2026-05-06-paths-explicit-and-Q3-refuted.tex}) gives an existence-and-explicit injection of paths with the right slack. The current theorem reduces the sign-positivity \emph{numerical} content to two single-component path-count bounds.

The two perspectives complement each other:
\begin{itemize}[leftmargin=2em,topsep=0.3em,itemsep=0.2em]
\item Theorem~B (path-level): asserts the existence of an injection of multisets of paths.
\item Theorem~\ref{thm:reduction}: identifies the load-bearing inequalities as $m^{(3,2)}_2 \geq 2$ and $m^{(4,1)}_1 \geq 2$.
\end{itemize}

A canonical path-level injection would simultaneously verify both inequalities via two specific injections of small finite sets, with the $\mathrm{atom}_{\mathrm{RTL}}$ slack identified as a specific extra path at $\bar c = 2$. Pinning down which specific path is the slack remains open (Section~\ref{sec:open}).

\section{What this does not give: canonical injection}\label{sec:open}

Theorem~\ref{thm:reduction} reduces sign-positivity to two integer inequalities. It does \emph{not} pin down the specific injection of multisets that achieves the inequality. The lex injection of \texttt{2026-05-06-paths-explicit-and-Q3-refuted.tex} is one specific injection; many others work equally well.

A canonical injection would identify, structurally, which path in $V_{(3,2)}$ at $\bar c = 2$ is ``the slack'' --- equivalently, which 2 of the 3 paths are ``imaged from $V_{(4,1)}$'' and which 1 is ``new'' to $V_{(3,2)}$.

We note one structural observation that does not, however, pin down the canonical injection:

\begin{remark}[Branching component of the slack at $\bar c = 2$]\label{rem:branching}
At $\bar c = 2$, the three paths in $V_{(3,2)}$ split by branching component as
\begin{itemize}[leftmargin=2em,topsep=0.2em,itemsep=0.1em]
\item Two paths with starting vertex in $V_{(3,1)}$-branching (vertex~$1$, lying in the $V_{(3,1)}$-summand of $V_{(3,2)} \downarrow S_4$).
\item One path with starting vertex in $V_{(2,2)}$-branching (vertex~$2$, lying in the $V_{(2,2)}$-summand).
\end{itemize}
Meanwhile, all five paths in $V_{(4,1)}$ at $\bar c = 2$ start at a $V_{(3,1)}$-branching vertex (vertex~$2$ of $V_{(4,1)}$, lying in the $V_{(3,1)}$-summand of $V_{(4,1)} \downarrow S_4$).

If the canonical injection respected the branching identification $V_{(3,2)} \downarrow V_{(3,1)} \;\leftrightarrow\; V_{(4,1)} \downarrow V_{(3,1)}$ on starting vertices, we would inject 5 $V_{(4,1)}$-paths into 2 (paths) $\times$ 2 (copies) $= 4$ slots --- one short. So at least one $V_{(4,1)}$-path must map to the $V_{(2,2)}$-branching slot.

This is a \emph{lower bound} on the non-canonical-by-branching content of the injection at $\bar c = 2$, but it does not single out a specific path.
\end{remark}

The corresponding observation at $\bar c = 0$ (where the slack is also 1) does not have the same flavor: the unique $V_{(3,2)}$ path at $\bar c = 0$ has starting vertex 0 (in the $V_{(3,1)}$-branching component of $V_{(3,2)}$), and the unique $V_{(4,1)}$ path at $\bar c = 0$ has starting vertex 0 (in the $V_{(4)}$-branching component of $V_{(4,1)}$). These are in different branching summands, but neither is forced to be the ``slack'' since the slack at $\bar c = 0$ is the $V_{(3,2)}$ path in copy~2 (i.e., a copy-choice rather than a path-choice).

\section{Open questions}

\begin{enumerate}[leftmargin=2em,topsep=0.3em,itemsep=0.3em]
\item \textbf{Structural lower bounds for the two multiplicity bounds.} Give a $W$-graph-natural proof that $m^{(3,2)}_2 \geq 2$ and $m^{(4,1)}_1 \geq 2$, without enumerating all closed paths.
\item \textbf{Canonical injection.} Pin down a $W$-graph-natural injection $\iota: \mathrm{Paths}(V_{(4,1)}) \hookrightarrow \mathrm{Paths}(V_{(3,2)}) \sqcup \mathrm{Paths}(V_{(3,2)})$ preserving the bigrading, with the slack identified as a specific structural feature.
\item \textbf{Generalization.} Does a similar reduction hold for other partition pairs $(\lambda, \mu)$ at $S_n$ with $A_\lambda - A_\mu$ in a similar simple form? Specifically, when $A_\lambda - A_\mu = N \cdot p(q)$ for some palindromic $p(q)$ and integer $N$, is the sign-positivity lemma $|N M_p| \leq M_\lambda$ reducible to a small number of multiplicity bounds, with the slack absorbed in a few specific components?
\end{enumerate}

\section*{Computational summary}

All computations performed in \texttt{\textasciitilde/projects/scratch/2026-05-06-sign-positivity/} and \texttt{\textasciitilde/projects/scratch/2026-05-06-paths/}. Cross-checks:
\begin{itemize}[leftmargin=2em,topsep=0.2em,itemsep=0.1em]
\item Multiplicity vectors $M_{(3,2)} = (1,5,3)$, $M_{(4,1)} = (1,3,5)$ (Apr~24, recomputed May~6).
\item $A_{(3,2)} - A_{(4,1)} = 2q(q^2+q+1)$ (Apr~29 polynomial residual zero).
\item $B_4$-decomposition $2q(q^2+q+1) = 0 + 2 q(1+q)^2 - 2 q^2$ (Lemma~\ref{lem:2q-decomp}).
\item Path counts $|\mathrm{Paths}(V_{(3,2)})_{\bar c}| = (1, 5, 3)$ and $|\mathrm{Paths}(V_{(4,1)})_{\bar c}| = (1, 3, 5)$ at $S_5$ (May~6, \texttt{path\_enumerate.py}).
\item Sign-positivity numerical: $\mathrm{atom}_{\mathrm{RTL}} = (1, 7, 1)$ and $\mathrm{atom}_{\mathrm{LTR}} = (1, 1, 7)$, both non-negative (May~6, \texttt{test\_theorem\_B.py}).
\end{itemize}

\end{document}
